\let\negmedspace\undefined
\let\negthickspace\undefined
\documentclass[journal]{IEEEtran}
\usepackage[a4paper, margin=10mm, onecolumn]{geometry}
%\usepackage{lmodern} % Ensure lmodern is loaded for pdflatex
\usepackage{tfrupee} % Include tfrupee package

\setlength{\headheight}{1cm} % Set the height of the header box
\setlength{\headsep}{0mm}  % Set the distance between the header box and the top of the text

\usepackage{gvv-book}
\usepackage{gvv}
\usepackage{cite}
\usepackage{amsmath,amssymb,amsfonts,amsthm}
\usepackage{algorithmic}
\usepackage{graphicx}
\usepackage{float}
\usepackage{textcomp}
\usepackage{xcolor}
\usepackage{txfonts}
\usepackage{listings}
\usepackage{enumitem}
\usepackage{mathtools}
\usepackage{gensymb}
\usepackage{comment}
\usepackage[breaklinks=true]{hyperref}
\usepackage{tkz-euclide} 
\usepackage{listings}
% \usepackage{gvv}                                        
\def\inputGnumericTable{}                                 
\usepackage[latin1]{inputenc}                                
\usepackage{color}                                            
\usepackage{array}                                            
\usepackage{longtable}                                       
\usepackage{calc}                                             
\usepackage{multirow}                                         
\usepackage{hhline}                                           
\usepackage{ifthen}                                           
\usepackage{lscape}
\usepackage{tikz}
\usetikzlibrary{patterns}

\begin{document}

\bibliographystyle{IEEEtran}
\vspace{3cm}

\title{12.82}
\author{ee25btech11063-vejith}

\maketitle
% \maketitle
% \newpage
% \bigskip
{\let\newpage\relax\maketitle}
\renewcommand{\thefigure}{\theenumi}
\renewcommand{\thetable}{\theenumi}
\setlength{\intextsep}{10pt} % Space between text and floats
\textbf{Question}:\\
let $\Vec{M}$ be a 3$\times$ 3 symmetric matrix with eigenvalues 1,2 and 3. let $\Vec{N}$ be any 3$\times$ 3  matrix with real eigenvalues such that $\Vec{MN} + \Vec{NM}= 3\Vec{I}$, where $\Vec{I}$ is a 3$\times$ 3  identity matrix. Then which of the following cannot be eigenvalue\brak{s} of the matrix $\Vec{N}$ ?
\hspace{2cm}\brak{\text{ST } 2022}
\begin{enumerate}[label=\alph*)]
    \begin{multicols}{4}
  \item $\frac{1}{4}$        
  \item $\frac{3}{4}$
  \item $\frac{1}{2}$
  \item $\frac{7}{4}$  
    \end{multicols}
\end{enumerate}

\textbf{Solution}:\\
let $\lambda_1$  , $\lambda_2$ , $\lambda_3$ be the eigenvalues of $\vec{N}$\\
By Eigenvalue decomposition
\begin{align}
    \vec{M}=\vec{QDQ}^\vec{\top}\\
     \vec{N}=\vec{Q}\vec{N}'\vec{Q}^\top
\end{align}
where 
\begin{align}
    \vec{Q}^\top\vec{Q}=\vec{I}\\
    \vec{D}=\text{diag}\brak{1,2 ,3}\\
    \vec{N}'=\text{diag}\brak{\lambda_1,\lambda_2 ,\lambda_3}\\
    \vec{MN}+\vec{NM}=3\vec{I}
\end{align}
From (1) and (2) (6) can be simplified as 
\begin{align}
        \vec{QD}\vec{N}'\vec{Q}^\top + \vec{Q}\vec{N}'\vec{DQ}=3\vec{I}\\
        \implies \vec{Q}^\top \brak{\vec{QD}\vec{N}'\vec{Q}^\top + \vec{Q}\vec{N}'\vec{DQ}}=\vec{Q}^\top\brak{3\vec{I}}\\
        \implies \brak{\vec{Q}^\top\vec{Q}}\vec{D}\vec{N}'\vec{Q}^\top + \brak{\vec{Q}^\top\vec{Q}}\vec{N}'\vec{DQ}=3\vec{Q}^\top\\
        \implies \vec{D}\vec{N}'\vec{Q}^\top + \vec{N}'\vec{DQ}=3\vec{Q}^\top\\
        \implies \brak{\vec{D}\vec{N}'\vec{Q}^\top + \vec{N}'\vec{DQ}}\vec{Q}= 3\vec{Q}^\top\vec{Q}\\
        \implies \vec{D}\vec{N}' + \vec{N}'\vec{D}=3\vec{I}
\end{align}
From (4) and (5)
\begin{align}
    \text{diag}\brak{\lambda_1,2\lambda_2 ,3\lambda_3} + \text{diag}\brak{\lambda_1,2\lambda_2 ,3\lambda_3} = 3\vec{I}\\
    \implies \text{diag}\brak{\lambda_1,2\lambda_2 ,3\lambda_3}=\frac{3}{2}\vec{I}\\
    \implies \begin{pmatrix}
        \lambda_1 & 0 & 0\\
        0 & 2\lambda_2 & 0\\
        0 & 0 & 3\lambda_3
    \end{pmatrix}=\begin{pmatrix}
        \frac{3}{2} & 0 & 0\\
        0 & \frac{3}{2} & 0\\
        0 & 0 & \frac{3}{2}
    \end{pmatrix}\\
    \implies \lambda_1=\frac{3}{2} \text{ and } \lambda_2=\frac{3}{4} \text{ and } \lambda_3=\frac{1}{2}
\end{align}
The eigen values of $\vec{N}$ are $\frac{3}{2}$ and $\frac{3}{4}$ and $\frac{1}{2}$
\end{document}